\section{Network-Aware Deep Learning and the Hybrid}
\label{sec:deep}

\subsection{The independent-segment baseline and its gap}
\label{sec:gap}
\textbf{XGBoost per segment.} Following the project brief, the classical baseline
of Phase~2 is a set of $N{=}207$ independent models, one per road segment. Model
$f_i$ sees only sensor $i$'s own history: 12 forward-filled lags; 1- and 3-step
differences; rolling mean, standard deviation and minimum; missing count and a
missing-now flag over the last hour; sin/cos time of day, day of week and a weekend
flag; the weekly profile $\mu_i(\tau(t{+}h))$ at all 12 target times (known in
advance), $\mu_i(\tau(t))$, and the current anomaly $a_{t,i}$ (37 features). The
brief also suggests flow and occupancy features; METR-LA publishes speed only, so
these would require the raw PeMS station files, which we leave as an extension.
Static road-segment attributes are constant within a per-segment model and
therefore carry no information there; the pooled gradient-boosting model of
Phase~1 uses them (Section~\ref{sec:features}). Each
$f_i$ is a multi-output XGBoost regressor~\cite{chen2016xgboost} with vector-leaf
trees (one tree predicts all 12 horizons), absolute-error loss, depth 6, learning
rate 0.1, and early stopping (patience 30) on every second validation window. On six
randomly chosen sensors, scored on the validation windows not used for early
stopping, vector-leaf trees were slightly more accurate than one tree per horizon
(MAE 2.721 vs.\ 2.738\,mph) at similar cost (3.4 vs.\ 2.9\,s per sensor;
\texttt{scripts/compare\_xgb\_strategy.py}).

\textbf{The gap, formally.} Let $\mathcal{F}^{i}_t$ be the information in sensor
$i$'s own past and $\mathcal{F}^{\mathcal{G}}_t\supseteq\mathcal{F}^{i}_t$ the
information in the whole network's past. By the law of total variance, the best
forecasts under squared loss, $m_i=\mathbb{E}[y\mid\mathcal{F}^{i}_t]$ and
$m_{\mathcal{G}}=\mathbb{E}[y\mid\mathcal{F}^{\mathcal{G}}_t]$, satisfy
\begin{equation}
\begin{split}
\mathbb{E}\big[(y-m_i)^2\big]-\mathbb{E}\big[(y-m_{\mathcal{G}})^2\big]
&=\mathbb{E}\big[(m_{\mathcal{G}}-m_i)^2\big]\ \ge 0,
\end{split}
\label{eq:gap}
\end{equation}
because $y-m_{\mathcal{G}}$ is orthogonal to every function of
$\mathcal{F}^{\mathcal{G}}_t\supseteq\mathcal{F}^{i}_t$. Equality holds only if the
network adds nothing. An independent model can at best reach the first term. The gap
is non-zero exactly when the residual
$r_{t+h,i}=y_{t+h,i}-f_i(\mathcal{F}^{i}_t)$ is predictable from other sensors at
time $t$. We test this directly on out-of-fold \emph{training} residuals, without
touching the test set: (i) are residuals of neighbouring segments correlated, and
(ii) does information available at time $t$, namely the anomalies and trends of
downstream and upstream neighbours, correlate with $r$? Section~\ref{sec:p2results}
shows that both hold. This is the empirical justification for adding a graph model.

\subsection{ST-GAT: spatio-temporal graph attention network}
\label{sec:stgat}
\textbf{Input.} For each window the network receives
$\mathcal{X}\in\R^{T'\times N\times 6}$ with $T'{=}12$. The channels are the
z-scored, causally forward-filled speed; a 0/1 observation mask; and $\sin/\cos$
of time of day and of day of week (Section~\ref{sec:dldata}). A learned node
embedding $\vect{e}_i\in\R^{16}$ is concatenated at every step~\cite{shao2022stid},
and a $1{\times}1$ convolution lifts the channels to $d{=}64$.

\textbf{Gated temporal convolution.} As in STGCN~\cite{yu2018stgcn}, a valid
convolution of width $K_t{=}3$ along time produces $2d$ channels split into $P,Q$:
\begin{equation}
\Gamma(\mathcal{H})=\big(P+\mathcal{R}(\mathcal{H})\big)\odot\sigma(Q),\qquad [P;Q]=\Theta\star_t\mathcal{H},
\label{eq:glu}
\end{equation}
where $\mathcal{R}$ crops the input to the output length (residual path) and the
sigmoid gate~\cite{dauphin2017glu} controls information flow.

\textbf{Directional graph attention.} The road graph is directed (73\,\% of
edges are one-way) and congestion moves against the direction of travel. We
therefore attend separately over \emph{downstream} neighbours
$\mathcal{N}^{\downarrow}(i)=\{j: W_{ij}>0\}\cup\{i\}$ and \emph{upstream}
neighbours $\mathcal{N}^{\uparrow}(i)=\{j: W_{ji}>0\}\cup\{i\}$, each with $H{=}4$
heads of width $d/(2H)$. For one direction and head,
\begin{align}
s_{ij}&=\operatorname{LeakyReLU}_{0.2}\!\big(\vect{a}_{\mathrm{dst}}^\top W\bar{\vect{h}}_i+\vect{a}_{\mathrm{src}}^\top W\bar{\vect{h}}_j\big),\\
\alpha_{ij}&=\frac{\exp(s_{ij})}{\sum_{k\in\mathcal{N}(i)}\exp(s_{ik})},\qquad
\vect{h}'_{t,i}=\sum_{j\in\mathcal{N}(i)}\alpha_{ij}\,W\vect{h}_{t,j},
\label{eq:gat}
\end{align}
which is exactly the GAT layer~\cite{velickovic2018gat}. A unit test checks our
implementation against PyTorch Geometric's \texttt{GATConv}~\cite{fey2019pyg} to
$10^{-5}$. One deliberate change is that the scores use the window-averaged state
$\bar{\vect{h}}_i=\frac{1}{T}\sum_t\vect{h}_{t,i}$. The coefficients are therefore
computed once per window and shared by its time steps, as in ASTGCN's spatial
attention~\cite{guo2019astgcn}. This turns each head into one masked $N{\times}N$
matrix product and needs $T$ times fewer attention scores ($HN^2$ instead of
$THN^2$) than per-step attention. The cost is that the weights cannot change
\emph{within} a one-hour window. The eight heads are
concatenated, passed through ELU, mixed by a linear layer, and added back
(residual).

\textbf{Block and output.} One ST block is the sandwich
$\mathcal{H}\leftarrow\operatorname{Dropout}\big(\operatorname{LN}(\Gamma_2(\mathcal{H}'+\operatorname{GAT}(\mathcal{H}')))\big)$
with $\mathcal{H}'=\Gamma_1(\mathcal{H})$ and layer normalisation~\cite{ba2016layernorm}.
Two blocks shrink the window $12{\to}10{\to}8{\to}6{\to}4$ steps, a final gated
convolution of width 4 collapses it to one vector per sensor, and a two-layer MLP
outputs all 12 horizons at once (direct strategy). The model has 163,516 parameters,
about half as many as Graph WaveNet (300,952).

\textbf{Receptive field and complexity.} Every output sees all 12 input steps.
Spatially, two blocks give two hops in each direction, which covers on average 36
sensors within two undirected hops; the median edge is 2.8\,km. One block costs
$O(BT'Nd^2K_t)$ for the convolutions and $O(BHN^2 + BT'HNd)$ for attention. With
$N{=}207$ the dense masked form is cheaper on a GPU than gathering the 1,515 edges.

\textbf{What attention adds.} The diffusion convolution of
Eq.~\eqref{eq:diffusion} uses fixed weights $(P^k)_{ij}$ that depend only on road
distance. A mean aggregator~\cite{hamilton2017sage} uses $1/|\mathcal{N}(i)|$.
Eq.~\eqref{eq:gat} makes $\alpha_{ij}$ a function of the current states, so the
influence of a neighbour can grow when it is congested. Setting
$\vect{a}_{\mathrm{src}}=\vect{a}_{\mathrm{dst}}=\vect{0}$ recovers the mean
aggregator, so attention strictly generalises it. Our ablation measures what this
adaptivity is worth. Attention, mean aggregation and the masked inputs are all
permutation-equivariant in the sensors; only the node embedding breaks the symmetry,
deliberately.

\subsection{Hybrid: XGBoost per segment \texorpdfstring{$\to$}{to} graph attention}
\label{sec:hybrid}
\textbf{Design.} The hybrid keeps what each model does best. The per-segment
XGBoost supplies a strong forecast of each segment's own level, seasonality and
recent trend $\hat{\vect{y}}^{\,\mathrm{seg}}_{i}\in\R^{12}$. ST-GAT learns only
what the independent models cannot: the network-level correction
\begin{equation}
\hat{\vect{y}}_{i}=\hat{\vect{y}}^{\,\mathrm{seg}}_{i}+g_\theta\big(\mathcal{X},\hat{Y}^{\mathrm{seg}};\mathcal{G}\big)_i .
\label{eq:hybrid}
\end{equation}
The base forecasts of all sensors enter $g_\theta$ through a linear embedding. They
are fused with the ST-GAT representation, and one more directional GAT layer lets
each sensor attend to its neighbours' forecasts. A segment can thus react when the
independent model just downstream predicts a slowdown. The last layer of $g_\theta$
is zero-initialised, so training starts \emph{exactly} at the XGBoost forecast (a
unit test checks this) and gradient descent only has to learn a residual. Residual
learning~\cite{he2016resnet,zhang2003hybrid} is easier: by
Section~\ref{sec:gap}, the residual has much lower variance than the speed itself,
and its structure is precisely the spatial propagation that graph attention models.

\textbf{Leakage-free stacking.} The second stage must be trained on base forecasts
that are as wrong as they will be at test time. In-sample forecasts are too good:
the XGBoost in-sample training MAE is 2.69\,mph against 2.99\,mph out-of-fold.
A meta-learner trained on them would learn to trust the base too much. We therefore
\emph{cross-fit}~\cite{wolpert1992stacking}. The training anchors are cut into
$K{=}5$ contiguous blocks. For each block an XGBoost model is fitted on the other
blocks, after removing anchors within 23 steps of the block so that no 23-step
window (11 input and 12 target steps) is shared, and it forecasts the held-out
block~\cite{bergmeir2012cv}. Validation and test anchors receive forecasts from
the model fitted on the whole training period, so evaluation stays strictly causal.
The ablation also trains the hybrid on in-sample base forecasts to show the cost of
skipping this step.

\subsection{Deep learning dataset and regularisation}
\label{sec:dldata}
\textbf{Dataset construction.} Samples are sliding 24-step windows (12 in, 12
out) indexed by their anchor time, split chronologically 70/10/20 with no shuffling
across the boundaries. Shuffling happens only \emph{within} the training period,
between mini-batches. The scaler, the weekly profile, the XGBoost models and every
statistic are fitted on training timestamps only.

\textbf{Outages are masked, not zero traffic.} The Phase~1 pipeline (inherited
from DCRNN) feeds a missing reading to the network as $0$\,mph, which is about
$-4.5$ standard deviations and indistinguishable from a total standstill. The new
pipeline forward-fills the last observation (causal) and adds an explicit
observation-mask channel~\cite{che2018grud}, so the network can learn how much to
trust a stale value. Missing \emph{targets} are excluded from the loss
(Eq.~\eqref{eq:loss}) and from every metric. An ablation trains ST-GAT on the old
zero-filled inputs.

\textbf{Cyclic time.} Time of day and day of week enter as
$(\sin 2\pi u,\cos 2\pi u)$. The Phase~1 scalar $u\in[0,1)$ places 23:55 and 00:00
at opposite ends of the scale, while the circle embedding keeps them adjacent and
keeps all slots at equal norm.

\textbf{Regularisation.} The networks have $1.6\times10^5$ parameters, and the
training set holds 23,974 overlapping windows that are strongly autocorrelated, so
the effective sample size is much smaller. We use dropout~\cite{srivastava2014dropout}
(0.2) after every block and in the head; attention dropout (0.1) on the
coefficients, which randomly removes neighbours, much like DropEdge; layer
normalisation; AdamW with decoupled weight decay $10^{-4}$~\cite{loshchilov2019adamw};
gradient clipping at norm 5; halving the learning rate after 3 epochs without
validation improvement; and early stopping~\cite{prechelt1998} with patience 8,
keeping the best validation checkpoint. Each epoch visits a fresh random half of the
training windows. Adjacent windows share 11 of 12 steps, so this halves the cost and
acts as mild data subsampling.
